\fsdivision{文件系统操作是一组状态转换协议}

前一节说明了 create 会修改位图、inode、目录项和数据块。结构清单还不足以构成正确
实现：并发线程可能同时修改同一目录，资源分配可能在中途失败，设备也可能只持久化
部分写入。因此，一个文件操作还需要规定加锁范围、错误回收、对外可见点和持久化顺序。

\fstopic{可见性、持久性与一致性}

分析操作结果时，需要分别观察三个边界：

\begin{description}
\item[可见性（visibility）] 其他线程现在能否通过名字或已打开的文件描述符观察到修改。
  它讨论的是运行中内存状态，不自动意味着数据已经掉电不失。
\item[持久性（durability）] 系统成功返回某个同步边界后，掉电重启是否仍能恢复该修改。
  它依赖文件系统提交协议、块层顺序和设备对 flush/FUA 的兑现。
\item[一致性（consistency）] 位图、inode、目录项和块映射是否仍满足格式规定的不变量。
  一致的文件系统也可能保存旧数据；“结构能挂载”不等于“应用的最后一次写已持久”。
\end{description}

\code{write} 返回通常建立的是新字节对进程的可见性；\code{rename} 可以建立
名字切换的原子可见性；\code{fsync} 请求建立持久性。三个边界有关联，却不是同一个
边界。把它们统称为“原子写”会掩盖真正需要证明的性质。

\fstopic{锁、事务、持久化顺序与回滚}

\emph{锁（lock）}是运行期间的互斥或同步规则。例如父目录锁可以防止两个 create
同时把同一个名字插入目录；inode 锁可以协调 truncate 与 write。锁保存在内存中，
断电后会消失，所以锁本身不能提供崩溃恢复。

\emph{事务（transaction）}把多项持久修改归为一个恢复单位。日志型文件系统先把
“本事务要发布哪些元数据版本”写入日志，并用 commit record（提交记录）表示事务
已经完整。恢复程序只重放具有有效提交证据的事务。事务解决的是“掉电后该接受哪组
修改”，并不替代运行中的锁。

\emph{持久化顺序（write ordering）}规定某个写在另一个写获得持久性之前必须先完成。
例如 ordered 模式需要先保证新文件数据已写出，再提交让该块可达的元数据，否则
重启后可能通过正确目录名读到旧块残留。设备可以重排普通写，因此文件系统还要通过
依赖关系、flush、FUA 或等价协议把逻辑顺序传到设备。

\emph{回滚（rollback）}处理的是尚未发布时的运行期错误。例如分配 inode 成功、分配
数据块却返回 ENOSPC，就要撤销 inode 位图和计数的修改。日志恢复常用的是重放已提交
事务，而不一定在磁盘上逐步执行传统意义的“反向撤销”。运行期回滚与崩溃恢复处理的
时间边界不同，所需证据也不同。

\begin{longtable}{@{}F{0.18\linewidth}F{0.36\linewidth}F{0.38\linewidth}@{}}
\toprule
机制 & 解决的问题 & 适用边界 \\
\midrule
锁 & 防止并发交错破坏运行状态 & 只在系统运行期间有效，不能跨越掉电 \\\tablerowrule
运行期回滚 & 回收系统调用中途失败后已取得的资源 & 在错误返回前恢复内存与事务状态 \\\tablerowrule
日志事务 & 识别一组元数据更新是否已经完整提交 & 用于 mount 恢复与 checkpoint；不自动提供普通数据内容原子性 \\\tablerowrule
持久化顺序 & 防止设备重排或易失缓存破坏写入依赖 & 约束 I/O 到达稳定存储的先后，不负责线程互斥 \\\tablerowrule
fsck 式扫描 & 修复已经违反全局不变量的磁盘结构 & 可恢复结构与空间，通常无法还原应用的原始意图 \\
\bottomrule
\end{longtable}

\fstopic{create：目录项是名字的可见性发布点}

设两个线程同时执行 \code{open("/d/x", O_CREAT|O_EXCL)}。如果都先查到名字不存在，
再分别插入，就可能产生重名目录项或泄漏其中一个 inode。因此检查和插入必须受同一
目录修改协议保护，不能在两步之间释放父目录的排他锁。

一次简化但完整的 create 可以按下面的状态机理解：

\begin{lstlisting}
create(parent, "x"):
  1. lock(parent directory)             // 串行化同目录名字修改
  2. lookup("x")                        // 在锁内再次确认不存在
  3. reserve journal credits            // 预留本事务最坏所需日志空间
  4. allocate inode; initialize it       // 还没有目录项，外部路径不可达
  5. if initial data is required:
       allocate block; write data
  6. insert dirent "x" -> inode          // 发布名字：路径开始可达
  7. update parent time/counts
  8. attach all metadata changes to transaction
  9. unlock(parent directory)

on error before step 6:
  free block if allocated
  free inode if allocated
  cancel or stop transaction
  unlock(parent directory)
  return error
\end{lstlisting}

第 6 步建立目录项，是新名字的可见性发布点。在此之前，新 inode 即使已经存在，也没有
从根目录出发的路径；在此之后，路径解析可以找到它。发布之前出错，资源应被回滚。
发布之后若向用户
返回失败，调用者也不能武断地认为名字不存在：网络文件系统、I/O 错误或崩溃边界都
可能造成“操作已经生效但答复丢失”，稳健程序要重新查询状态。

日志空间预留发生在结构修改之前。这样，即使操作走到最坏路径，也有足够空间记录维持
一致性所需的元数据；若等目录项已经发布后才发现日志已满，提交和回滚都会失去可靠
空间保障。

\fstopic{create 各阶段的掉电结果}

下面假设没有日志，只按顺序写 MiniFS-Lab 的几个块。表中的“泄漏”是指位图认为资源
已占用，却没有从根目录可达的对象拥有它。

\begin{longtable}{@{}F{0.24\linewidth}F{0.28\linewidth}F{0.38\linewidth}@{}}
\toprule
掉电点 & 重启后可能看到 & 形成原因 \\
\midrule
只写 inode 位图后 & 一个已占用但未初始化/不可达的 inode & 分配状态先于对象正文持久 \\\tablerowrule
再写块位图后 & 不可达 inode，加一个无所有者数据块 & 两种资源都分了，但映射尚未建立 \\\tablerowrule
再写 inode 与数据后 & 完整数据仍没有名字可达 & 目录项尚未持久 \\\tablerowrule
只持久目录项、inode 未持久 & 名字指向无效或旧 inode & 发布发生在被发布对象之前 \\\tablerowrule
全部写完后 & 结构可能正确 & 仍要确认设备是否兑现写序与持久化命令 \\
\bottomrule
\end{longtable}

日志把相关元数据版本连同事务号写入日志，完整写出 commit
record 后才把事务视为可重放；没有完整提交证据的尾部事务被忽略。ordered data 模式
还建立“数据先于使其可达的元数据提交”的顺序。恢复时由提交记录区分完整事务与尾部
半成品：前者重放，后者忽略。

\fstopic{unlink：名字消失后的对象生命周期}

目录项对 inode 的持久引用数由链接计数表示；进程已打开文件所持有的是运行时引用。
这两种引用必须分开。考虑：

\begin{lstlisting}
fd = open("/tmp/report", O_RDWR);  // 得到运行时 file/inode 引用
unlink("/tmp/report");             // 删除名字，链接计数可能变为 0
write(fd, "tail", 4);              // 仍然合法
close(fd);                          // 最后一个运行时引用消失
\end{lstlisting}

\code{unlink} 的直接语义是从父目录移除“名字 $\rightarrow$ inode”映射并减少链接计数。
链接计数降为零后，新的路径查找找不到该对象；但只要 \code{fd} 仍引用它，内核就不能
回收 inode 和数据块。最后一个打开引用关闭后，回收路径才释放块映射和 inode 编号。

这里产生一个新的崩溃窗口：目录项已删除、链接计数为零，但系统在释放全部数据块前
掉电。很多日志型文件系统维护类似 orphan list（孤儿链表）的持久恢复线索。它记录
“这个 inode 已经没有名字，但清理尚未完成”；mount 恢复时继续 truncate/删除，避免
永久泄漏。这条记录是删除协议主动保存的清理进度，不是恢复程序根据残留块进行猜测。

\fstopic{rename：原子可见性与持久发布是两回事}

常见安全更新模式是先写临时文件，再把它改名到正式名字：

\begin{lstlisting}
fd = open("config.tmp", O_CREAT | O_TRUNC | O_WRONLY, 0600);
write_all(fd, new_bytes);
fsync(fd);                         // 先持久化临时文件内容
close(fd);
rename("config.tmp", "config");  // 名字切换对并发查找原子可见
dirfd = open(".", O_DIRECTORY | O_RDONLY);
fsync(dirfd);                      // 请求持久化目录项变化
close(dirfd);
\end{lstlisting}

这里的\emph{原子可见}表示其他进程不会观察到半个目录项：目标名在某个线性化点从旧
对象切到新对象。但 \code{rename} 成功返回本身不应被泛化成“掉电后新名字必然存在”。
文件数据和目录项是两个持久对象，所以先同步文件，再执行 rename，再同步父目录。
跨两个目录 rename 时，若应用要求两个目录的变更都可靠，还要按目标平台契约处理
相关目录同步。

rename 还暴露了锁顺序问题。两个线程若一个先锁目录 A 再锁 B，另一个先锁 B 再锁 A，
就会形成死锁：双方都拿着对方需要的锁等待。VFS 因此规定目录操作的锁层级与祖先
关系规则；跨目录 rename 还必须防止把目录移动进自己的后代，从而制造目录环。由此，
锁顺序同时承担并发安全和目录无环不变量的维护职责。

\fstopic{truncate：同时收缩长度、缓存和块映射}

\code{truncate(file, 5000)} 看似只改一个长度字段，实际上旧文件若长 20 KiB，它还要：

\begin{enumerate}\tightlist
\item 在 inode 锁与映射失效协议保护下发布新的 \code{i\_size}；
\item 处理 page cache 中新 EOF 之后的 folio，不能让旧缓存内容再次被写回；
\item 将包含新 EOF 的最后一个块尾部清零，避免文件以后增长时泄露旧字节；
\item 删除 EOF 之后的 extent/间接块映射并释放物理块；
\item 与正在进行的 buffered write、direct I/O、writeback 和 mmap page fault 协调；
\item 把长度、块映射和空闲空间修改纳入可恢复的元数据事务。
\end{enumerate}

顺序错误会形成真实漏洞。如果先把物理块分给文件 B，再允许文件 A 的旧脏 folio 写回，
A 的写回就可能覆盖 B 的数据；如果不清理最后一个部分块，A 随后扩展时可能读到截断前
的旧内容。完整的 truncate 需要撤销 EOF 之后各层次对旧数据的引用，并与仍可能使用
这些引用的并发路径建立同步屏障；修改 \code{i\_size} 只是其中一步。

对已 mmap 的文件，映射地址也不是永久拥有物理块。截短后访问新 EOF 之外的页面可能
触发 \code{SIGBUS}；内核需要让相关页表映射和 page cache 状态失效。这个例子再次说明：
磁盘 inode、内存 inode、page cache 与进程页表是同一文件的不同层次，不能只改其中一个。

\fstopic{fsync 的对象范围与调用顺序}

对普通文件调用 \code{fsync(fd)}，目标是该文件的数据和恢复它所需的相关元数据；
\code{fdatasync} 可以不强求与数据读取无关的部分元数据，但仍必须同步影响数据正确
读取的长度等信息。二者都必须检查返回值，因为延迟分配、writeback 或设备错误可能
直到同步时才报告。

文件名属于目录数据，不属于普通文件 inode。因此新建文件的可靠发布通常还需要同步
父目录；rename、link、unlink 改的也是目录命名关系。不能因为文件内容已同步，就推出
重启后包含它的名字也一定存在。反过来，只同步目录也不能替代先把新文件内容同步好。

\textbf{持久性检查应明确对象集合和顺序。}需要同步哪些文件数据、inode 元数据与父
目录项，各操作采用什么先后顺序，返回错误怎样处理，以及设备是否提供正确的持久化
语义。应用若需要跨多个文件原子提交，通常还要使用 WAL、manifest 或数据库事务；
单个文件的 \code{fsync} 不会自动把多个文件合并成一个事务。

本节涉及的操作可以用四条不变量统一检查：

\begin{enumerate}\tightlist
\item 每个可达目录项都指向一个已分配、类型合法的 inode；
\item 每个可达 inode 引用的数据块都已分配，且非共享格式中不会被别的 inode 重复拥有；
\item 不可达且无运行时引用的零链接 inode 最终能够回收，清理中断后仍有恢复线索；
\item 对外声称已经持久的状态，必须存在一条从有效超级块/日志提交点到该状态的完整恢复路径。
\end{enumerate}

前两条是空间与引用一致性，第三条是生命周期闭环，第四条是崩溃后的可证明性。后面的
page cache、日志和 COW 看似是新主题，其实都在实现或优化这四条约束。
